\documentclass[12pt]{article}

\usepackage[margin=1in]{geometry}
\usepackage[hidelinks]{hyperref}

\begin{document}

\noindent
Dear Professor Kirillova and Professor Masiero,\\

\noindent
Please consider my manuscript, \textit{``WOW Factor: Measuring Labor Loss from a Negative Tourism Supply Shock in Iceland,''} for publication in \textit{Annals of Tourism Research Empirical Insights}. This manuscript has not been previously published and is not currently under consideration by any other journal. I confirm that I have no competing interests to disclose.\\

\noindent
The paper examines the labor market effects of the unexpected 2019 bankruptcy of WOW Air, a major low-cost airline in Iceland, and uses this event as a natural experiment to identify the consequences of a sudden tourism supply shock. International connectivity is essential for tourism-dependent economies, particularly island destinations, yet it also creates a critical vulnerability. Despite broad recognition of this dependence, much of the existing literature relies on descriptive vulnerability indexes rather than clean causal evidence of how these disruptions matriculate.\\

\noindent
Using a Bayesian Structural Time Series model, I estimate counterfactual outcomes for Iceland in the absence of the bankruptcy and find that WOW Air’s collapse reduced passenger departures by more than 20\%, tourism employment by over 8\%, and national unemployment by nearly 50\% relative to the counterfactual path. The findings demonstrate both the economic value from affordable air connectivity and the substantial labor market consequences when that connectivity is abruptly lost.\\


\noindent
These results have direct policy relevance. Governments are often hesitant to intervene when a major carrier faces financial distress, and it can be difficult to justify assisting a failing private organization or softening the effects of bankruptcy. This paper provides evidence that some firms occupy uniquely important positions within the tourism supply chain, where collapse can create economy-wide labor market consequences. The results also identify which sectors and workers are most exposed, helping inform where targeted mitigation and recovery efforts may be most effective.\\

\noindent
I believe this manuscript would be of strong interest to the readers of \textit{Annals of Tourism Research Empirical Insights} because it contributes causal evidence to the literature on tourism vulnerability, resilience, and labor market adjustment while offering clear implications for destination management and policy design.\\


\noindent
Thank you for your time and consideration.\\


\noindent
Sincerely,

\vspace{1em}

\noindent
Luke F. Miller \\
Ph.D. Candidate, Department of Economics \\
University of Hawaiʻi at Mānoa \\
Email: \href{mailto:lukefm@hawaii.edu}{lukefm@hawaii.edu} \\
Phone: (412) 807-8570 \\
ORCID: \href{https://orcid.org/0009-0008-2932-1568}{0009-0008-2932-1568}

\end{document}